\uuid{fRGg}
\exo7id{4218}
\titre{exo7 4218}
\auteur{quercia}
\organisation{exo7}
\datecreate{2010-03-11}
\isIndication{false}
\isCorrection{true}
\chapitre{Equation différentielle}
\sousChapitre{Equations aux dérivées partielles}
\cpge{mpsi-mp2i.primitives-equations-differentielles.c, pcsi.techniques-calcul-primitives.c, ptsi.techniques-calcul-primitives.c, mpsi-mp2i.fonctions-deux-variables.b, pcsi.fonctions-deux-variables.b, ptsi.fonctions-deux-variables.b, mpsi-mp2i.fonctions-deux-variables.c, pcsi.fonctions-deux-variables.c, ptsi.fonctions-deux-variables.c, bcpst1.derivees-primitives-integrales.b, bcpst1.equations-differentielles-lineaires.b, bcpst1.fonctions-deux-variables.c, mp-mpi.calcul-differentiel-optimisation.b, pc.fonctions-plusieurs-variables.c, psi.fonctions-plusieurs-variables.c}
\module{Analyse}
\niveau{L2}
\difficulte{}

\contenu{
\question{
Soit $f:\R \to \R$ de classe $\mathcal{C}^2$ et $g(x,y,z) = \frac {f(r)}r$
avec $r = \sqrt{x^2+y^2+z^2}$.
Déterminer $f$ de sorte que $\frac{\partial^2 g}{\partial x^2} + \frac{\partial^2 g}{\partial y^2} + \frac{\partial^2 g}{\partial z^2} = -g$.
}
\reponse{
$f(r) = A\cos r + B\sin r$.
}
}
